\section{兼容历史与 x86-64 架构状态}

x86-64 不是从 64 位模式重新设计的一套孤立处理器。8086 的分段、IBM PC 的
I/O 端口、80286 的保护机制、80386 的 32 位扩展和 AMD64 的长模式依次叠加，
复位仍保留早期入口，现代内核必须沿这条兼容路径主动建立新状态。

\topic{一、这里说的“历史”是在解释今天的接口}
处理器历史可以回答一个现代代码问题：为什么一颗 64 位 CPU 复位后不直接
进入 64 位，为什么还有 16 位段寄存器、\code{IN/OUT}、A20、GDT、PIC
和 PS/2？这些机制保留下来，因为新机器仍要运行大量已经存在的软件和扩展卡。

\term{兼容性}{compatibility}至少有三层，不能混成“新机器还能开机”：

\begin{osgridlongtable}{L{0.20\linewidth}L{0.31\linewidth}L{0.39\linewidth}}
层次 & 问题 & x86 例子 \\ \hline
源代码兼容 & 旧源码是否只需少量修改就能重新编译 & 许多 C 接口能在 32/64 位重新编译，但指针宽度变化 \\ \hline
二进制兼容 & 已经生成的机器码能否直接执行 & x86-64 长模式保留兼容子模式运行许多 32 位程序 \\ \hline
平台兼容 & 固件、地址、端口、设备和中断约定是否仍存在 & PC 仍保留传统 COM、PIC、PIT 等兼容接口或模拟 \\ \hline
\end{osgridlongtable}

处理器 ISA 兼容只保证指令和架构状态，不能自动保证某个操作系统能驱动所有
主板。IBM PC 兼容平台把 CPU 之外的固件入口、I/O 端口和扩展总线约定也固定
下来，才形成“同一份软件能在不同厂商机器上运行”的生态。

\topic{为什么厂商愿意背负旧接口}%
已经部署的软件、开发工具、用户数据和扩展卡具有巨大价值。若新处理器只能运行
全新软件，用户换机就要同时替换整个生态；保留兼容入口能让新硬件先运行旧系统，
再由新操作系统逐步启用新能力。代价是新设计不能自由删除旧状态，启动过程会
出现多层模式切换。

\topic{兼容不等于所有历史器件都还焊在板上}%
现代 SoC 不需要放置一颗与 1980 年完全相同的独立 8259A 芯片。它可以在片上
逻辑、芯片组或虚拟机里实现相同的软件可见寄存器；也可以只在固件早期提供兼容
窗口，操作系统随后切换到 APIC、PCIe 和现代定时器。对软件而言，关键是当前
接口契约是否存在，而不是硅片内部是否沿用旧版晶体管布局。

\topic{2.1 8086：16 位计算如何看到 1 MiB 地址空间}
1978 年的 8086 以 16 位为主要寄存器宽度。一个 16 位无符号数最多表达
\(2^{16}=65536\) 种值，只能直接编号 64 KiB 位置；而处理器有 20 根地址线，
可以表达 \(2^{20}=1\) MiB 物理字节地址。若每个地址都必须放进单个 16 位
寄存器，硬件能力便会浪费。

8086 采用“段值左移四位再加偏移”：
\[
\mathrm{physical}=(\mathrm{segment}\ll4)+\mathrm{offset}.
\]
段寄存器给出 16 字节粒度的基址，16 位偏移在其上选择最多 64 KiB 窗口。
例如 \code{0x1000:0x0020} 形成
\[
\code{0x1000}\times16+\code{0x0020}=\code{0x10020}.
\]
这在有限晶体管和寄存器宽度下扩大了可寻址范围，也允许程序通过调整段值搬动
代码或数据窗口。

代价是地址不唯一。\code{0x1000:0x0020} 与
\code{0x1001:0x0010} 指向同一个物理字节；跨过 64 KiB 边界时还要处理段。
20 位总线只保留低 20 位，最大和可能回绕到低地址。后来 PC 为兼容这种回绕
行为引出 A20 gate；这条历史会在长模式启动章节再次出现。

\topic{2.2 8088 与 IBM PC：CPU 之外的平台约定变成生态}
8088 与 8086 使用相近指令架构，但外部数据总线较窄，便于当时用较便宜的
8 位板级器件构建系统。IBM 在 1981 年的 PC 中选择 8088，并公开了足够多的
平台结构，使兼容机厂商可以围绕相同软件和扩展卡生产机器。

这里真正长寿的不只是 CPU 指令，还包括一整套主板选择：

\begin{itemize}[leftmargin=2em]
  \item 固件在特定高地址提供复位入口和 BIOS 服务；
  \item 8259A 管理外部中断，8253/8254 提供定时；
  \item 键盘控制器、串口和磁盘控制器占用约定 I/O 端口；
  \item ISA 总线让扩展卡在约定地址、IRQ 和 DMA 通道工作；
  \item DOS 和应用程序直接依赖这些接口，而非抽象设备树。
\end{itemize}

早期硬件资源很少，固定约定能减少发现和配置成本；设备增多以后，固定 IRQ 与
端口又产生冲突，最终促成 Plug and Play、PCI 枚举和 ACPI。现代 PC 同时带着
两种历史：启动早期仍提供兼容入口，成熟操作系统则尽快发现并使用现代资源。

\topic{2.3 80286：多任务需要保护，不只是更快运算}
实模式程序可以写任意可寻址内存，一个错误指针便能破坏操作系统或其他程序。
多用户、多任务系统需要回答：某段内存属于谁，普通程序是否可以改中断表，
越界时 CPU 能否在写入发生前阻止？

80286 引入保护模式。段寄存器不再直接充当段基址，而是成为
\term{选择子}{descriptor selector}；CPU 用它索引 GDT/LDT 描述符，取得基址、
界限、类型和特权级。一次内存访问同时接受地址计算和权限检查，违规会进入异常，
而不是依赖程序自觉。

这种方案解决了隔离，但 286 保护模式和实模式切换不够灵活，旧 DOS 软件又大量
存在。平台和操作系统需要复杂兼容策略。这个经验说明：新增安全状态并不意味着
旧入口消失，模式间可控切换本身也是架构能力。

\topic{2.4 80386：32 位寄存器与分页让地址空间更规则}
80386 把通用寄存器扩到 32 位，EAX、EBX、ESP、EIP 等能直接表达 4 GiB
范围。保护模式继续存在，并增加分页：线性地址先由页目录和页表翻译为物理页，
每页还带 present、读写和用户权限。

分页让操作系统不必要求每个进程的物理内存连续。两个进程可以各自认为程序位于
同一线性地址，页表却把它们送到不同物理页；暂时不用的页还可换出到磁盘。段
提供较粗粒度的基址、界限与权限，页表提供固定页粒度映射。后来的 32 位系统
常把段基址设为零、界限覆盖整个线性空间，把主要隔离交给分页。

80386 仍从兼容复位状态启动。新能力不是电源一开就自动生效：固件或内核必须
准备描述符、页表和异常入口，再设置控制位进入保护/分页状态。这正是 STUOS
逐级模式切换的历史来源。

\topic{2.5 80486、Pentium 与现代微架构：接口稳定，内部越来越并行}
随后处理器把浮点、缓存和更深流水线集成进 CPU，Pentium 时代又发展超标量、
分支预测、乱序执行和更复杂缓存。软件并不会因为内部同一时刻执行多条微操作，
就突然观察到任意乱序结果；架构规定异常、内存顺序和指令提交边界。

这一时期还逐步加入 CPUID、APIC、PAE、MMX/SSE、MSR 和多处理器能力。每项
扩展都需要发现机制和启用顺序。操作系统不能只因为“CPU 看起来很新”就执行
某条扩展指令；它先用 CPUID 检查能力，再配置保存现场和控制寄存器。本项目
v1.2/v1.3 对 FXSR、SSE2 与 SYSCALL 的规格门禁就是同一原则。

\topic{2.6 AMD64：64 位能力仍从旧入口进入}
到 32 位系统后期，4 GiB 虚拟地址空间限制大型数据库、文件缓存和内核布局。
PAE 扩大物理地址，却没有把单个普通进程的线性指针自然扩成 64 位。AMD64
扩展通用寄存器到 64 位，增加 R8--R15，引入四级页表和长模式；Intel 后来
实现兼容的 Intel 64，x86-64 生态由此形成。

长模式内部又区分 64 位子模式和兼容子模式。64 位代码获得 RIP 相对寻址、更宽
寄存器和新的调用 ABI；兼容子模式帮助运行旧 32 位代码。大部分传统分段在
64 位子模式中被弱化，但 \code{FS}/\code{GS}、权限检查、描述符入口和复位
兼容仍有作用。

最关键的遗留是：x86-64 CPU 复位后仍不会直接成为分页开启的 64 位处理器。
软件要从兼容入口建立 A20、GDT、PAE 页表、CR3、EFER.LME 和 CR0.PG，再通过
远控制转移装入 64 位代码段。这条链看似绕远，实际上是“旧机器可启动”与
“新系统可用 64 位”同时成立的兼容契约。

\begin{osgridlongtable}{L{0.15\linewidth}L{0.22\linewidth}L{0.27\linewidth}L{0.26\linewidth}}
代际 & 代表能力 & 软件获得的新边界 & 保留下来的兼容条件 \\ \hline
8086 & 16 位寄存器、20 位地址 & 分段实模式与 1 MiB 空间 & 复位入口、段寄存器 \\ \hline
80286 & 保护模式、描述符 & 权限、界限与任务机制 & 仍从实模式复位 \\ \hline
80386 & 32 位寄存器、分页 & 4 GiB 线性地址、页级保护 & 16/32 位模式并存 \\ \hline
Pentium 系列 & cache、乱序、APIC & 性能和多处理器能力 & 指令集二进制兼容 \\ \hline
AMD64 & 64 位寄存器、四级页表 & 长模式和更大地址空间 & 继续从兼容复位状态进入 \\ \hline
\end{osgridlongtable}

\topic{怎样使用这张代际表}%
表中“新边界”表示软件从该代开始能够依赖的能力，“兼容条件”表示新能力没有
删除的旧入口。它不是说每台现代机器都把所有旧外设原样焊在主板上；QEMU 可以
模拟接口，现代 SoC 可以集成兼容逻辑，固件也可以只在启动早期使用后交给 APIC
和 PCIe。

\section{寄存器、标志与模式控制状态}

架构状态是软件可通过规范观察的寄存器、标志、内存和异常结果。通用寄存器、
RIP、RFLAGS、段寄存器、控制寄存器和 MSR 属于这类状态。流水线深度、重排序
缓冲、物理寄存器数量和分支预测表属于微架构实现；它们影响性能，却不能改变
正确程序最终观察到的指令语义。

\term{ISA}{Instruction Set Architecture，指令集架构}是二进制软件与 CPU
之间的主要契约。它规定某串机器码怎样解码、操作数多宽、哪些寄存器改变、
何时产生异常。两颗处理器可以都实现 x86-64 ISA，却采用不同缓存大小、执行
单元和分支预测器；同一 STUOS 镜像在架构能力满足时仍应得到相同语义结果。

\term{微架构}{microarchitecture}回答“这款 CPU 内部怎样高效实现 ISA”。
现代 CPU 可能先把 x86 指令拆成内部微操作，预测跳转，提前执行后续操作，再
按程序语义提交。若预测错误，未提交结果必须撤销；若指令在架构上产生精确异常，
操作系统看到的 RIP 和寄存器现场必须落在规范定义的位置。

\begin{osgridlongtable}{L{0.20\linewidth}L{0.31\linewidth}L{0.39\linewidth}}
状态类别 & 软件怎样观察 & 为什么操作系统关心 \\ \hline
通用寄存器 & 指令读写、异常现场保存 & 函数参数、返回值、地址和线程切换 \\ \hline
控制寄存器/MSR & 特权指令读写 & 模式、分页、CPU-local、系统调用入口 \\ \hline
内存与设备状态 & load/store、MMIO、Port I/O & 代码数据、页表、驱动协议 \\ \hline
异常与中断 & CPU 改变控制流并压入/保存现场 & 故障隔离、时钟、设备完成 \\ \hline
缓存/预测器 & 通常不能作为普通架构值直接依赖 & 影响性能、一致性和侧信道，不改变合法结果 \\ \hline
\end{osgridlongtable}

\topic{“不可见”不表示完全没有影响}%
缓存会改变访问延迟，多核缓存一致性影响并发可见顺序，预测器还可能形成侧信道。
这些机制不能让正确执行的加法返回另一个数，却会影响性能、安全和内存模型。
初学阶段先建立架构正确性，后续再在明确边界上处理微架构效应。

\begin{pdfdiagram}
  \node[os state] (software) {软件可见\\ISA 契约};
  \node[os block,below left=of software] (regs) {寄存器\\RAX--R15、RIP};
  \node[os block,below=of software] (control) {控制状态\\CR0、CR3、CR4、EFER};
  \node[os block,below right=of software] (memory) {内存与 I/O\\地址、数据、异常};
  \node[os hardware,below=22mm of control] (micro) {微架构\\流水线、cache、预测、执行单元};
  \draw[os arrow] (software) -- (regs);
  \draw[os arrow] (software) -- (control);
  \draw[os arrow] (software) -- (memory);
  \draw[os arrow] (regs) -- (micro);
  \draw[os arrow] (control) -- (micro);
  \draw[os arrow] (memory) -- (micro);
\end{pdfdiagram}
\diagramnote{ISA 规定软件可见结果，微架构负责在物理资源上实现这些结果。
操作系统依赖前者，性能调优和侧信道分析才进一步关心后者。}

\topic{四、通用寄存器的宽度为什么一层层叠加}
x86-64 的 RAX 低 32 位叫 EAX，低 16 位叫 AX，AX 的低高字节分别叫 AL 和
AH。四个名称指向同一个物理架构值的不同视图。写 EAX 会把 RAX
高 32 位清零；写 AX 或 AL 只更新相应低位。固件使用 AL 访问 8 位 UART，
使用 AX 配置 16 位段寄存器和栈，原因来自硬件接口与当前模式的精确宽度。

名字也保存了历史。8086 有 AX、BX、CX、DX、SP、BP、SI、DI 八个主要 16 位
寄存器；80386 在前面加 E 形成 32 位 EAX、ESP 等；AMD64 再加 R 形成 RAX、
RSP，并增加 R8--R15。新名字扩展旧值而不是删除旧视图，所以一条 8 位设备指令
仍可使用 AL，一条 64 位地址计算可使用 RAX。

\begin{osgridlongtable}{L{0.16\linewidth}L{0.19\linewidth}L{0.25\linewidth}L{0.30\linewidth}}
宽度 & RAX 中的名字 & 一次写入的结果 & 常见用途 \\ \hline
64 位 & \code{RAX} & 覆盖全部 64 位 & x86-64 地址、整数、系统调用结果 \\ \hline
32 位 & \code{EAX} & 覆盖低 32 位并把高 32 位清零 & 32 位运算、用零扩展构造 64 位值 \\ \hline
16 位 & \code{AX} & 只覆盖低 16 位 & 早期 x86 状态、16 位端口或字段 \\ \hline
低 8 位 & \code{AL} & 只覆盖 bits 0--7 & 8 位 UART 数据和字符 \\ \hline
高 8 位 & \code{AH} & 只覆盖 bits 8--15 & 历史接口；与某些新编码组合受限 \\ \hline
\end{osgridlongtable}

“CPU 是 64 位”不代表每次访问都应使用 64 位。设备寄存器、磁盘字段和 ABI
指定多宽就必须访问多宽；对 8 位 UART 做 64 位输出可能触碰不存在或相邻寄存器。
反过来，保存物理地址的内部量若只用 32 位，会在大内存机器上截断。

\begin{center}
\begin{tikzpicture}[os diagram,x=0.65cm]
  \draw[BookBlue,thick] (0,0) rectangle (16,1);
  \draw[BookTeal,thick] (8,0) rectangle (16,1);
  \draw[BookAmber,thick] (12,0) rectangle (16,1);
  \draw[BookRed,thick] (14,0) rectangle (16,1);
  \node at (4,0.5) {RAX 高 32 位};
  \node at (10,0.5) {EAX 高 16 位};
  \node at (13,0.5) {AH};
  \node at (15,0.5) {AL};
  \node[below] at (8,-0.1) {RAX 64 位};
  \node[below] at (12,-0.7) {EAX 32 位};
  \node[below] at (14,-1.3) {AX 16 位};
\end{tikzpicture}
\end{center}

\topic{五、段寄存器为什么同时有可见值和隐藏缓存}
\term{段寄存器}{segment register}包括 \code{CS}、\code{DS}、\code{SS}、
\code{ES}、\code{FS} 和 \code{GS}。字母最初分别暗示代码、数据、栈和附加
数据用途。程序能够直接看到的是一个 16 位值，但处理器内部还为每个段寄存器
保存一组软件不能像普通寄存器那样直接读写的缓存：段基址、界限、类型、特权
属性和是否存在等。文档常把前者称为\term{可见部分}{visible part}，把后者称为
\term{隐藏部分}{hidden descriptor cache}。

为什么需要隐藏缓存？保护模式中的描述符包含的信息远多于 16 位，若处理器每次
取指或访存都重新访问 GDT/LDT，会增加多次内存读取和权限检查。于是加载段
寄存器时完成一次描述符查找与合法性检查，把结果缓存下来；后续地址形成直接
使用缓存。描述符表后来被修改，并不会自动刷新已经加载的段缓存，必须重新加载
相应段寄存器。

\begin{osgridlongtable}{L{0.16\linewidth}L{0.24\linewidth}L{0.25\linewidth}L{0.25\linewidth}}
处理器状态 & 可见值表示什么 & 隐藏缓存怎样得到 & 本项目要关心什么 \\ \hline
普通实模式 & 直接参与 16 字节粒度分段 & 加载时把段值左移 4 位形成基址 & 16 位入口地址和 64 KiB 窗口 \\ \hline
复位瞬间 & \code{CS} 看起来像历史实模式值 & CPU 预置特殊高基址 & 第一条指令位于高地址复位向量 \\ \hline
保护模式 & GDT/LDT 的描述符选择子 & CPU 读取描述符的基址、界限和权限 & 代码段、数据段和特权检查 \\ \hline
64 位子模式 & 多数段基址/界限被弱化 & \code{CS} 属性仍决定 64 位代码；\code{FS}/\code{GS} 基址仍有用 & 长模式入口、CPU-local 和线程数据 \\ \hline
\end{osgridlongtable}

复位状态正好展示了“可见值不等于完整状态”。复位后 \code{CS} 的隐藏基址为
\code{0xFFFF0000}，指令偏移为 \code{0xFFF0}，相加得到
\code{0xFFFFFFF0}。这就是 32 位地址空间顶部附近的复位向量。第一次正常重载
\code{CS} 会按当前模式重新建立隐藏缓存，特殊复位基址随之消失。

\topic{为什么切换代码模式需要远跳转}%
\term{近跳转}{near jump}只改变当前代码段中的指令偏移，通常不重新加载
\code{CS}；\term{远跳转}{far jump}同时给出新的代码段选择子和偏移。设置
\code{CR0.PE} 或完成长模式控制位配置后，CPU 的控制位虽然已经改变，当前
\code{CS} 缓存却仍可能描述旧代码环境。远跳转让处理器检查新的代码段描述符、
刷新 \code{CS} 隐藏缓存，并从明确的新入口继续执行。

\code{iret}/\code{iretq} 也能恢复比普通 \code{ret} 更多的架构状态，例如
\code{CS}、\code{RIP}、\code{RFLAGS}，跨特权级时还恢复栈状态。它们不是
“更长的函数返回”，而是异常、硬件事件或特权转换协议的一部分。若栈上的字段
顺序、选择子或权限不合法，返回动作本身也会产生异常。

\topic{六、RFLAGS：同一运算为什么有多种“溢出”}
RFLAGS 是一组一位状态的集合，而不是一个普通整数。算术指令会同时写入多种
结果解释：进位 CF 面向无符号数，溢出 OF 面向有符号数，零 ZF 表示结果全零，
符号 SF 复制结果最高位。CPU 并不知道一个位模式在源代码中究竟叫
\code{unsigned} 还是 \code{signed}；它产生两套信息，后续条件跳转选择正确
的一套。

以 8 位加法为例：

\begin{itemize}[leftmargin=2em]
  \item \code{0xFF + 0x01 = 0x00}（只保留低 8 位）时，CF 为 1、ZF 为 1、
        OF 为 0。无符号的 \(255+1\) 超出 8 位，但有符号的
        \(-1+1=0\) 没有溢出。
  \item \code{0x7F + 0x01 = 0x80} 时，CF 为 0、OF 为 1、SF 为 1。
        无符号的 \(127+1=128\) 可以表达；有符号 8 位最大只有 127，
        两个正数却得到负号位，因此发生有符号溢出。
\end{itemize}

所以“检查溢出”必须先回答正在解释有符号数还是无符号数。汇编中的无符号条件
常查看 CF/ZF，有符号条件还组合 SF/OF/ZF。若中间插入另一条会改标志的算术或
\code{test} 指令，前一次比较结果就被覆盖；比较与条件跳转通常要紧邻安排。

中断允许 IF 和方向 DF 不描述算术，而是控制处理器行为。入口执行 \code{cli}
清 IF，避免在 IDT 和栈尚未建立时接收可屏蔽中断；执行 \code{cld} 清 DF，
让 \code{lodsb}、\code{movsb} 等字符串指令向高地址推进。

\begin{osgridlongtable}{L{0.12\linewidth}L{0.23\linewidth}L{0.55\linewidth}}
标志 & 产生或修改方式 & 固件中的意义 \\ \hline
CF & 算术、\code{stc}/\code{clc} & 无符号进位；也可作为固件函数的一位返回协议 \\ \hline
ZF & 比较、测试和算术结果 & \code{test al,al} 判断字符串结尾 \\ \hline
SF/OF & 算术结果的最高位与有符号溢出 & 正确解释有符号比较和范围错误 \\ \hline
IF & \code{sti}/\code{cli} & IDT 建立前保持可屏蔽中断关闭 \\ \hline
DF & \code{std}/\code{cld} & 控制字符串指令地址增加或减少 \\ \hline
\end{osgridlongtable}

\topic{IF 不是整个中断系统的总电源}%
一个键盘按键最终成为 CPU 可屏蔽中断，至少要经过四层许可：设备自身开启事件，
中断控制器没有屏蔽该输入，控制器把事件路由到目标 CPU，CPU 的 IF 为 1。
任一层关闭，普通 IRQ 都不会按预期进入处理函数。反过来，\term{异常}{exception}
由正在执行的指令同步产生，\term{不可屏蔽中断}{NMI} 有独立规则，它们不会因为
\code{cli} 清除 IF 就全部消失。因此“已经关中断”只能准确表述为“CPU 暂不接受
可屏蔽外部中断”。

\topic{七、控制寄存器如何把同一颗 CPU 推进到新模式}
CR0 包含保护模式与分页总开关；CR3 指向当前页表根；CR4 选择 PAE 等扩展；
EFER 是 MSR，其中 LME 允许进入长模式。这些位不是互相独立的按钮。长模式
激活需要 PAE 页表、CR3、EFER.LME 与 CR0.PG 按规范顺序共同成立。

\term{控制寄存器}{control register}保存会改变整颗逻辑处理器解释指令和地址
方式的状态，只有高特权代码可以修改。\term{MSR}{Model-Specific Register，
模型专用寄存器}是另一类用编号访问的控制状态；EFER 位于 MSR 空间。名称中的
“模型专用”并不表示可以随意猜编号，而是必须依据 CPUID 与处理器规范确认该
实现支持哪些寄存器和位。

从复位环境进入 64 位代码，概念依赖链如下。实际实现可因固件入口状态而跳过
已经满足的步骤，但绝不能假设依赖不存在。

\begin{osgridlongtable}{L{0.07\linewidth}L{0.22\linewidth}L{0.30\linewidth}L{0.31\linewidth}}
序号 & 状态或动作 & 它解决的问题 & 做错以后常见现象 \\ \hline
1 & 关闭可屏蔽中断并建立已知栈 & 防止半初始化现场被异步打断 & 随机跳转、栈破坏 \\ \hline
2 & 确认 A20 地址线可用 & 防止 1 MiB 以上地址按旧兼容规则回绕 & 页表或代码写到低地址 \\ \hline
3 & 构造 GDT 并执行 \code{lgdt} & 提供合法的保护/64 位代码段描述符 & 加载选择子时保护异常 \\ \hline
4 & 设置 \code{CR0.PE} 并远跳转 & 进入保护模式并刷新代码段缓存 & 继续按错误位宽解码指令 \\ \hline
5 & 构造 PML4 等页表并设置 \code{CR4.PAE} & 提供长模式要求的分页格式 & 开分页时异常 \\ \hline
6 & 把页表物理根地址写入 \code{CR3} & 告诉硬件从哪里开始页表遍历 & 取指映射不存在 \\ \hline
7 & 通过 \code{wrmsr} 设置 \code{EFER.LME} & 声明分页开启后允许长模式 & 仍停留在非长模式状态 \\ \hline
8 & 设置 \code{CR0.PG} & 启用地址翻译并激活长模式条件 & 下一次取指即可能缺页 \\ \hline
9 & 远跳转到 64 位代码段与已映射入口 & 刷新 \code{CS}，正式按 64 位代码执行 & 错误指令解码或保护异常 \\ \hline
\end{osgridlongtable}

这不是九个彼此独立的“开关”。例如，写入 \code{CR3} 的是页表根的\emph{物理}
地址；页表本身以及模式切换后的下一条指令都必须位于当前映射可达范围；
\code{CR4.PAE}、\code{EFER.LME} 和 \code{CR0.PG} 必须形成规范要求的组合。
每次修改后，CPU 会立刻按新规则解释接下来的取指或访存，不会等函数结束后再
统一生效。

\topic{为什么模式切换错误常表现成“机器突然重启”}%
若开启分页后下一条指令没有映射，CPU 产生缺页异常；若 IDT 尚未有效，或者异常
处理入口本身也无法取指，CPU 无法交付第一次异常并尝试双重故障；双重故障仍
无法交付时进入三重故障，平台通常让处理器复位。屏幕来不及显示 C++ 错误信息，
QEMU 窗口看起来只像一闪或重启。因而模式切换代码要短，每跨过一个危险边界就
输出可观测标记，并预先检查下一条指令、栈、GDT 和页表是否可达。

复位阶段保持这些复杂机制关闭，只建立最小可观察环境。后续每个里程碑只推进
一组状态，使异常发生时仍能定位最后成功边界。

\topic{八、特权、系统调用和异常为什么必须由硬件约束}
保护模式引入特权级和描述符检查，分页继续增加页级读写执行权限。非法指令、
除零、缺页和保护错误不会返回普通函数错误码，而是由处理器切换控制流到异常
向量。内核必须在执行可能失败的机制之前建立 IDT 和诊断路径。

\term{特权级}{privilege level}是 CPU 每次取指、访问页和执行敏感指令时
都会检查的硬件状态，与桌面界面中的“管理员账户”含义不同。x86 定义 Ring 0 到
Ring 3 四级，现代通用操作系统主要使用 Ring 0 运行内核、Ring 3 运行用户
程序。当前特权级 CPL 来自代码段状态；描述符的 DPL、选择子的 RPL、页表的
U/S 位和读写位共同约束访问。普通程序即使把一段字节伪装成 \code{mov cr3}
也不能因此获得内核权限，CPU 会拒绝在 Ring 3 执行该特权操作。

\topic{用户程序怎样合法请求内核服务}%
用户程序不能直接改页表或驱动磁盘，但它必须能够打印、分配内存和读文件。
\term{系统调用}{system call}就是受控入口：用户代码按照 ABI 把系统调用号和
参数放入约定寄存器，执行 \code{syscall} 或其他规定入口；CPU 按预先配置的
控制状态进入内核代码，内核保存现场、检查编号和参数、执行被允许的服务，再
通过匹配的返回协议回到用户态。

硬件完成“从哪个入口进入、哪些架构状态改变”的一部分工作，却不会自动证明
用户指针安全。内核仍须检查地址是否属于允许的用户范围、长度加法是否溢出、
页面是否存在、访问期间是否会发生竞态。特权转换只建立边界，输入验证才让这个
边界真正安全。

\begin{osgridlongtable}{L{0.16\linewidth}L{0.23\linewidth}L{0.25\linewidth}L{0.26\linewidth}}
事件 & 谁发起 & 与当前指令的关系 & 典型用途或原因 \\ \hline
系统调用 & 用户程序主动执行约定指令 & 同步、可预期 & 请求内核服务 \\ \hline
异常 & CPU 检测当前指令或访问出错 & 同步；返回点由异常类型规定 & 除零、非法指令、缺页、保护错误 \\ \hline
IRQ & 外部设备经中断控制器发出 & 异步，可被 IF 和控制器屏蔽 & 时钟、键盘、网卡完成通知 \\ \hline
NMI & 平台通过不可屏蔽通道发出 & 异步，不受普通 IF 屏蔽 & 严重硬件事件或监控 \\ \hline
\end{osgridlongtable}

复位固件尚未建立异常处理，错误指令可能连续触发异常并最终三重故障复位。
串口标记和逐阶段实现的价值正在这里：最后一条已出现标记界定了故障所在区间。

\topic{九、把四十多年的历史缩回本项目的一条启动链}
历史名词只有落到一条可执行链上才真正有用。把前面的年代重新排列为 STUOS
开机时必须完成的因果关系，可以得到下面这条路径：

\begin{enumerate}[leftmargin=2.4em]
  \item \textbf{8086 留下复位与实模式入口。}CPU 从兼容复位向量取第一条
        指令，段和偏移决定早期地址；这解释了为什么第一阶段不是普通
        x86-64 C++ 函数。
  \item \textbf{IBM PC 留下平台兼容接口。}传统串口、键盘控制器、PIC 和
        固件约定让极早期代码可以用很少驱动建立观测；在 QEMU 中，这些设备
        由模型提供，在现代实体 SoC 中可能由片上兼容逻辑提供。
  \item \textbf{80286 留下描述符和特权概念。}启动代码准备 GDT，选择合法
        代码段，并让 CPU 开始按保护规则检查访问，而不是允许任意实模式程序
        改写所有内存。
  \item \textbf{80386 留下 32 位控制状态与分页。}页表把程序使用的线性地址
        翻译为物理地址，CR0/CR3/CR4 形成后续长模式入口的基础。
  \item \textbf{AMD64 留下长模式与新寄存器。}软件设置 PAE、EFER.LME 和
        分页条件，再通过远跳转进入 64 位代码，随后才能使用本书的
        x86-64 C++ 内核 ABI。
  \item \textbf{现代 PC 又增加 APIC、PCIe、ACPI、UEFI 和多核。}它们不是
        把前五步全部抹掉，而是在兼容入口之后提供可枚举设备、电源描述、
        高速总线和现代固件服务。QEMU 教学路径可以先固定最小模型，实体载板
        路径则必须读取固件和总线描述，不能把虚拟机参数当成铜线事实。
\end{enumerate}

因此，一次现代 x86-64 启动并不是“CPU 先假装成古董，再突然变现代”这么模糊。
更准确的描述是：处理器从最小、兼容、可预测的架构状态开始，软件逐项提供新
模式所需的数据结构，逐项打开能力，并在每个边界承担验证责任。旧机制解释
入口为何存在，新机制解释目标怎样工作，模式切换代码则是两者之间可审计的桥。
